\documentclass{article}

\usepackage{tgschola}
\usepackage[parfill]{parskip}
\usepackage{hyperref}

\begin{document}
	
	\section{Introduction}
	
	This is the documentation for the Tabda package. It is structured chronologically
	according to how you would run a debating tournament. My hope is that this makes 
	it easy to find the relevant information as you tab your tournament.
	
	Tabda is open source 
	software\footnote{Meaning that the source code is available to you, and may be re-
	used with very minimal restrictions. The source code can be found online
	\url{https://github.com/teymour-aldridge/tabda}. Tabda is made available subject
	to the terms of the BSD 3-Clause license.} maintained by Teymour Aldridge, with
	contributions from the community. Please report bugs to our issue tracker.
	
	\tableofcontents
	
	\section{Importing data}
	
	The first step in any tournament is to import user data. There are two ways to do
	this:
	
	\begin{enumerate}
		\item manually, from the participants page
		\item by uploading a CSV\footnote{Comma-separated values. Both Excel and
		Google Sheets provide a way to easily export data to a CSV.} file
	\end{enumerate}
	
	Support for team registration is planned, but not currently implemented.
	
	\subsection{Importing using the participants page}
	
	\subsection{Uploading from a CSV file}
	
	
	\section{The group stage}
	
	\section{The break}
	
	Computing the break is an important stage in Tabbycat.
	
	\section{The elimination stage}
	
	Elimination rounds
	
\end{document}